\documentclass[10pt,a4paper]{amsart}
\usepackage[margin=19mm]{geometry}
\usepackage{amsmath,amssymb,mathtools}
\usepackage[scr=rsfs,cal=euler]{mathalfa}
\usepackage{xfrac}
\usepackage[unicode,hidelinks,bookmarks=false]{hyperref}
\newcommand{\ie}{i.e.\ }
\newcommand{\cf}{cf.\ }
\newcommand{\resp}{respectively\ }
\newcommand{\Subsection}{\subsection}
\providecommand{\nobreakdash}{\nobreak\hbox{-}}
\setlength{\emergencystretch}{2em}
\multlinegap0pt
\usepackage{xcolor}
\usepackage{amssymb}
\usepackage{algorithm}\usepackage{algpseudocode}
\newtheorem{prop}{Proposition}
\title{RANDSMAPs: Random-Feature/multi-Scale Neural Decoders with Mass Preservation}
\author{Dimitrios G. Patsatzis, Alessandro Della Pia, Lucia Russo, Constantinos Siettos}
\date{Source: arXiv:2601.14794v1 (CC BY 4.0)}
\begin{document}
\maketitle

\noindent\emph{Source and licence.} RANDSMAPs: Random-Feature/multi-Scale Neural Decoders with Mass Preservation. Version arXiv:2601.14794v1 (CC BY 4.0), distributed by arXiv under the Creative Commons Attribution 4.0 International licence, http://creativecommons.org/licenses/by/4.0/, as recorded in the arXiv licence metadata for this exact version and shown on the abstract page https://arxiv.org/abs/2601.14794v1. Full licence notice: \texttt{licenses/arxiv-2601.14794-CC-BY-4.0.txt}.\par\smallskip

\noindent\textit{Editorial scope.} Counted unit: the article's prop prp:RANDSMAP (source0.tex lines 770--781). The article's complete original proof of it is delivered with it (source0.tex lines 782--862, one \texttt{proof} environment of 81 lines). No mathematics is written, repaired or re-derived: the statement and the proof are the article's own lines, in the article's own notation, and no retained line is substituted or re-flowed.  Retained uncounted as the source-local context the proof needs: lines 194--197; lines 347--351; lines 373--377; lines 396--399.  Nothing was omitted from any retained span, so the package carries no omission record.  No retained line contains a citation, so the wrapper carries no bibliography and no bibliography entry.  No retained line contains a figure environment, an image-inclusion command or drawing code, so no image is delivered and the package declares no asset dependency.  Editorial formatting only, in addition: an AMS \texttt{amsart} preamble that restates the article's own declarations and macros used by the retained lines, namely the environments \texttt{prop} on separate counters, together with the standard packages those lines need.  The retained environments print in this document as Proposition 1 in order of appearance, whereas the article prints the same items with the numbering of its own full document.  The complete native source of the article as distributed in the arXiv e-print for this exact version is bundled unchanged as \texttt{sources/193107/source0.tex}.
\begin{equation}
    1_M^\top \widehat{X} = 1_N^\top,
    \label{eq:s2one_cons_recon}
\end{equation}

\begin{equation}
\Phi_S = \big[1_n \, | \, \widetilde{\Phi}_S \big]
\in\mathbb{R}^{n\times (P+1)}, \quad \widetilde{\Phi}_{S_{ik}}=\phi(y_{s_i}; c_{k}),
\label{eq:RFNNBasis}
\end{equation} 

\begin{equation}\label{eq:obj-row}
\min_{A_S\in\mathbb{R}^{P+1\times M}} \; 
\big\|X^{\top}_S - \Phi_S A_S \big\|_2^2 + \lambda\,\|A_S\|_2^2,
\qquad \lambda>0,
\end{equation}

\begin{equation}
A_S=V_r\,(\Sigma_r^{2}+\lambda I_r)^{-1}\,\Sigma_r\, U_r^{\top} X_S^{\top},
\label{eq:TikhSVD}
\end{equation}

\begin{prop} \label{prp:RANDSMAP}
Suppose the (down)sampled dataset $\mathcal{X}_S=\{x_{s_i}\}_{i=1}^n\subset\mathbb{R}^M$ is mass-preserving, i.e., the data matrix $X_S \in \mathbb{R}^{M \times n}$ satisfies the (normalized) sum-to-one constraint $1_M^\top X_S = 1_n^\top$. Let $\Phi_S = U_r\Sigma_r V_r^\top$ be the SVD of the feature matrix $\Phi_S\in\mathbb{R}^{n\times (P+1)}$ with $\operatorname{rank}(\Phi_S)=r$. Then, the coefficient matrix $A_S\in\mathbb{R}^{(P+1)\times M}$ given by 
\begin{equation}
A_S=V_r(\Sigma^2_r+\lambda I_r)^{-1}\Sigma_r U^{\top}_r\big(X^{\top}_S-\frac{1}{M}\big[I_n-U_r(I_r + \lambda \Sigma^{-2}_r)U^{\top}_r\big]1_n\,1_M^\top\big),
    \label{eq:ShurAfinal0}
\end{equation}
or, in the vanishing regularization limit $\lambda \to 0$, by the Moore-Penrose pseudo-inverse-based solution
\begin{equation} A_S=V_r\Sigma^{-1}_rU^{\top}_rX^{\top}_S=\Phi_S^{+}X^{\top}_S,
    \label{eq:ShurAfinal0s}
\end{equation}
ensures that the reconstructed fields $\widehat{X}_S = A^{\top}_S \Phi^\top$ provided by RANDSMAP decoder are mass-preserving, satisfying the sum-to-one constraint $1_M^\top\widehat{X}_S=1_n^\top$.
\end{prop}

\begin{proof}
When the (down)sampled data satisfy mass conservation, the RANDSMAP reconstruction $\widehat{X}_S \in \mathbb{R}^{M \times n}$ must preserve it as in Eq.~\eqref{eq:s2one_cons_recon}, i.e., 
\begin{equation}
1_M^\top \widehat{X}_S = 1_n^\top \Leftrightarrow 1_M^\top A_S^\top \Phi_S^\top  = 1_n^\top \Leftrightarrow \Phi_S A_S\,1_M = 1_n.%\quad \text{or} \quad \frac{1}{n}1^{\top}_n \, K_S A_S\,1_M = 1,
\label{eq:masscond}
\end{equation}
To enforce this constraint, we will first derive the solution in Eq.~\eqref{eq:ShurAfinal0} for the regularized least‑squares problem in Eq.~\eqref{eq:obj-row} subject to the linear condition in Eq.~\eqref{eq:masscond}. The minimum $L_2$-norm solution in Eq.~\eqref{eq:ShurAfinal0s} is obtained in the vanishing-limit of the Tikhonov regularization solution.

Introducing a Lagrange multiplier $\mu\in\mathbb{R}^n$ for the $n$ equality constraints yields the Lagrangian
\begin{equation}
\mathcal{L}(A_S,\mu) = \frac{1}{2}\|X^{\top}_S-\Phi_S A_S\|^2_2 + \frac{\lambda}{2}\|A_S\|^2_2 + \mu^\top (\Phi_S A_S\,1_M -1_n), \, \lambda>0,
\end{equation}
which is minimized w.r.t. $A_S$ and $\mu$. Assuming the standard convexity hypothesis of $\nabla^2_{A_S A_S}\mathcal{L}(A,\mu)=\Phi_S^\top \Phi_S + \lambda I_{P+1}$ being positive definite, which is satisfied under the mild condition $\lambda>0$, and setting the derivatives of $\mathcal{L}(A_S,\mu)$ w.r.t. $A_S$ and $\mu$ to zero gives the system  
\begin{equation}
\label{eq:KKT-matrix}
\begin{aligned}
&(\Phi_S^\top \Phi_S + \lambda I_{P+1}) A_S + \Phi_S^\top \mu\,1_M^\top = \Phi_S^\top X^{\top}_S,\\
& \Phi_S A_S\,1_M =1_n.
\end{aligned}
\end{equation}
Equation \eqref{eq:KKT-matrix} is a linear system coupling unknowns $A_S\in\mathbb{R}^{(P+1)\times M}$ and $\mu\in\mathbb{R}^n$, which can be solved using the Shur complement, as $\Phi_S^\top \Phi_S + \lambda I_{P+1}$ is invertible. Solving first for $A_S$ implies
\begin{equation}
    A_S=(\Phi_S^\top \Phi_S + \lambda I_{P+1})^{-1}(\Phi_S^\top X^{\top}_S-\Phi_S^\top \mu\,1_M^\top),
    \label{eq:ShurA}
\end{equation}
which, upon substitution into $\Phi_S A_S\,1_M =1_n$, yields the reduced equation
\begin{align}
    \Phi_S(K_S^\top \Phi_S + \lambda I_{P+1})^{-1}(\Phi_S^\top X^{\top}_S-\Phi_S^\top \mu\,1_M^\top)1_M & = 1_n, \, \mbox{or} \nonumber \\
    \Phi_S(\Phi_S^\top \Phi_S + \lambda I_{P+1})^{-1}(\Phi_S^\top 1_n-M\,\Phi_S^\top \mu) & = 1_n, \, \mbox{or} \nonumber \\
    \Phi_S(\Phi_S^\top \Phi_S + \lambda I_{P+1})^{-1}\Phi_S^\top( 1_n-M\,\mu) & =1_n,
   \label{eq:shur3}
\end{align}
where the (transposed) constrain $X_S^\top 1_M=1_n$ was used in the second step. Using the dual form identity
\begin{equation}
    (\Phi_S^\top \Phi_S + \lambda I_{P+1})^{-1}\Phi_S^\top=\Phi_S^\top(\Phi_S \Phi_S^\top + \lambda I_n)^{-1},
    \label{eq:dual_identity}
\end{equation}
the expression in Eq.~\eqref{eq:shur3} becomes
\begin{equation}
\Phi_S\Phi_S^\top(\Phi_S \Phi_S^\top + \lambda I_n)^{-1}( 1_n-M\,\mu) =1_n,
   \label{eq:shur4}
\end{equation}

Now considering the SVD of the feature matrix $\Phi_S=U_r\Sigma_r V^{\top}_r$, where $U_r \in \mathbb{R}^{n\times r}$, $\Sigma_r\in \mathbb{R}^{r \times r}$ and $V_r \in \mathbb{R}^{(P+1)\times r}$ with $r=\operatorname{rank}(\Phi_S)$, substitution in Eq.~\eqref{eq:shur4} yields the solution:
\begin{equation}
    \mu =\frac{1}{M}\big[I_n-U_r(I_r + \lambda \Sigma_r^{-2})U_r^{\top}\big]1_n.
    \label{eq:eqmu}
\end{equation}
Finally, substituting the expression for $\mu$ in Eq.~\eqref{eq:eqmu} into the Shur component for $A_S$ in Eq.~\eqref{eq:ShurA} and using the SVD of $\Phi_S$, we obtain the closed-form mass preserving solution
\begin{equation}
A_S=V_r(\Sigma^2_r+\lambda I_r)^{-1}\Sigma_r U^{\top}_r\big(X^{\top}_S-\frac{1}{M}\big[I_n-U_r(I_r + \lambda \Sigma^{-2}_r)U^{\top}_r\big]1_n\,1_M^\top\big),
    \label{eq:ShurAfinal}
\end{equation}
recovering the expression in Eq.~\eqref{eq:ShurAfinal0}. We note that at the limit $M\to \infty$, the above solution simplifies to the Tikhonov spectral filter solution in Eq.~\eqref{eq:TikhSVD} for the unconstrained problem.

Let us now verify that the solution for $A_S$ in Eq.~\eqref{eq:ShurAfinal} is indeed mass preserving, by evaluating the sum-to-one constraint $\Phi_S A_S\,1_M=1_n$ in Eq.~\eqref{eq:masscond}. Starting from the left-hand side, substitution of $A_S$ and the SVD of $\Phi_S$ yields
\begin{align}
    \Phi_S A_S1_M = U_r\Sigma_rV_r^{\top }V_r(\Sigma^2_r+\lambda I_r)^{-1}\Sigma_r U^{\top}_r\big(X^{\top}_S-\frac{1}{M}\big[I_n-U_r(I_r + \lambda \Sigma^{-2}_r)U^{\top}_r\big]1_n\,1_M^\top\big)1_M &= \nonumber \\ 
    U_r\Sigma_r(\Sigma^2_r+\lambda I_r)^{-1}\Sigma_r U^{\top}_r\big(X^{\top}_S\, 1_M-\big[I_n-U_r(I_r + \lambda \Sigma^{-2}_r)U^{\top}_r\big]1_n\big) &= \nonumber\\ 
    U_r\Sigma_r(\Sigma^2_r+\lambda I_r)^{-1}\Sigma_r U^{\top}_r\big(1_n-\big[I_n-U_r(I_r + \lambda \Sigma^{-2}_r)U^{\top}_r\big]1_n\big) &= \nonumber \\ 
    %U_r\Sigma_r(\Sigma^2_r+\lambda I_r)^{-1}\Sigma_r U^{\top}_r\big[I_n-I_n+U_r(I_r + \lambda \Sigma^{-2}_r)U^{\top}_r\big]1_n=\\ \nonumber
     U_r\Sigma_r(\Sigma^2_r+\lambda I_r)^{-1}\Sigma_r U^{\top}_r\big(U_r(I_r + \lambda \Sigma^{-2}_r)U^{\top}_r\big)1_n &= \nonumber\\ 
     U_r\Sigma_r(\Sigma^2_r+\lambda I_r)^{-1}\Sigma_r (I_r + \lambda \Sigma^{-2}_r)U^{\top}_r\,1_n=U_r\, U^{\top}_r\,1_n, &
     \label{eq:ShurAconst}
\end{align}
where the (transposed) constrain $X_S^\top 1_M=1_n$ was used in the third step. Now, recall that due to the presence of output biases, the feature matrix is written as $\Phi_S=\big[1_n \, | \, \widetilde{\Phi}_S \big]$; see Eq.~\eqref{eq:RFNNBasis}. Making the trivial assumption that $1_n$ does not, in general, belong in the column space of $\widetilde{\Phi}_S$, i.e., $1_n \notin \mathcal{R}(\widetilde{\Phi}_S)$, the inclusion of output bias implies 
\[\Phi_S \, e_{bias}=1_n.
\] 
where $e_{bias}^\top=[1,0,\ldots,0]\in \mathbb{R}^{P+1}$; therefore implying that $1_n$ belongs in the column space of $\Phi_S$, i.e., $1_n \in \mathcal{R}(\Phi_S)$. But $U_r$ spans $\mathcal{R}(\Phi_S)$ and thus, 
\begin{equation}
U_rU^{\top}_r\,1_n=1_n.
\label{eq:UrUrt}
\end{equation}
Substitution of the above to Eq.~\eqref{eq:ShurAconst} results in $\Phi_S A_S\,1_M=1_n$, so that the condition for the preservation of the the mass is satisfied for the reconstructions $\widehat{X}_S$.

We finally note that the spectral form solution in Eq.~\eqref{eq:ShurAfinal} recovers, in the vanishing regularization limit $\lambda\to0$, the Moore-Penrose pseudo-inverse based solution
\begin{equation}
A_S=V_r\Sigma^{-1}_rU^{\top}_r\,X^{\top}_S=\Phi_S^{+}\, X^{\top},
\end{equation}
which also satisfies the mass preservation condition, due to the result in Eq.~\eqref{eq:UrUrt}.
\end{proof}

\end{document}
